% GAIOL Poster — Engineering Applications of Artificial Intelligence
% Compile: pdflatex GAIOL_poster.tex
% Place paper figures in ./figures/ (Figure_2.png).

\documentclass[14pt]{article}
\usepackage[utf8]{inputenc}
\usepackage[T1]{fontenc}
\usepackage{amsmath,amssymb}
\usepackage{graphicx}
\usepackage{booktabs}
\usepackage{multicol}
\usepackage{caption}
\usepackage[hidelinks]{hyperref}
\usepackage[
  paperwidth=36in,
  paperheight=24in,
  margin=1in,
  noheadfoot
]{geometry}
\usepackage{array}
\usepackage{xcolor}
\usepackage{titlesec}

\setlength{\columnsep}{1.5em}
\setlength{\parindent}{0pt}
\setlength{\parskip}{0.5em}
\setlength{\columnseprule}{0pt}

\definecolor{sectionbg}{RGB}{240,248,255}
\definecolor{accent}{RGB}{0,51,102}

\titleformat{\section}{\large\bfseries\color{accent}}{}{0em}{}[\vspace{-0.3em}\titlerule\vspace{0.4em}]

\begin{document}
\pagestyle{empty}

% ----- Title block -----
\begin{center}
\vspace{0.4cm}
{\Huge\bfseries Global Artificial Intelligence Operating Layer}\\[0.35em]
{\LARGE\bfseries Adaptive Bayesian Consensus for Multi-Model LLM Orchestration}
\vspace{0.5cm}

{\Large Ch. Sai Sathvik, D.\,V.\,S. Monish Kumar, Abhishek Vinod\\
Ramdas Kapila$^{1}$, Sumalatha Saleti$^{2}$}
\vspace{0.25cm}

{\large
$^{1}$Dept. of CSE (Data Science), B.V. Raju Institute of Information Technology, Telangana, India\\
$^{2}$Dept. of CSE, SRM University AP, Telangana, India}
\vspace{0.25cm}

{\normalsize \textit{Engineering Applications of Artificial Intelligence}}
\vspace{0.5cm}
\end{center}

\sloppy
\begin{multicols}{3}

% ----- Abstract -----
\section*{Abstract}
Modern AI deployments rely on multiple LLMs and agents but are orchestrated via ad-hoc scripts without principled model selection or cross-provider consensus. We present the \textbf{Global Artificial Intelligence Operating Layer (GAIOL)}, a layered orchestration framework that coordinates heterogeneous LLM providers behind a uniform interface, decomposes complex queries into parallelizable subtasks, and aggregates responses through a novel \textbf{Adaptive Bayesian Trust-Weighted Consensus (ABTC)} mechanism. ABTC maintains per-model, per-domain Beta-distributed trust variables updated online after each consensus round. On a 500-query benchmark (analytical reasoning, code generation, multi-step problem solving, knowledge retrieval, creative synthesis), GAIOL achieves overall quality $0.83 \pm 0.02$ (24\% above single-model, 13\% above LangChain), 95.2\% success rate, and 5\,ms orchestration overhead. An ablation study confirms ABTC yields statistically significant gains over uniform and hand-tuned static consensus ($p < 0.01$) across all five domains.

% ----- Problem \& contributions -----
\section*{Problem and Contributions}
Today's ecosystem offers no satisfactory orchestration layer: models are tied to vendor APIs, agent frameworks use incompatible protocols, and orchestration logic lives in brittle pipelines. We propose GAIOL and evaluate:
\begin{itemize}
\item \textbf{Layered orchestration architecture}: provider-agnostic interface, Universal AI Protocol, adapter-based integration; provisions for federated data and governance.
\item \textbf{Global orchestrator}: beam-search task decomposition, strategy-based model routing, multi-model response assembly.
\item \textbf{ABTC algorithm}: per-model, per-domain trust learned online via Beta-Bernoulli updates, replacing static consensus weights.
\end{itemize}

% ----- System architecture -----
\section*{System Architecture}
GAIOL follows a layered microservices design (see figure): Presentation (web, API gateway); Business Logic (reasoning engine, orchestration, RAG, consensus); Data Access (multi-tenant storage, vector search); External Integration (Gemini, Hugging Face, OpenRouter adapters). The reasoning engine implements a multi-stage pipeline: decomposition $\rightarrow$ model selection $\rightarrow$ parallel inference $\rightarrow$ scoring $\rightarrow$ ABTC consensus $\rightarrow$ verification.

\begin{center}
\IfFileExists{figures/Figure_2.png}{%
  \includegraphics[width=\columnwidth]{figures/Figure_2}%
  \captionof{figure}{System architecture}}{%
  \fbox{\parbox{0.85\columnwidth}{\centering\small\textit{Add figures/Figure\_2.png (system architecture) from paper.}}}}
\end{center}

% ----- Algorithmic framework -----
\section*{Algorithmic Framework}
\textbf{Five phases:} (1) Decomposition: $Q \rightarrow T = \{t_1,\ldots,t_n\}$; (2) Smart orchestration: map $t_i$ to $\mathcal{M} \subset \mathcal{R}$; (3) Heterogeneous inference: parallel candidates $C_i$; (4) Consensus: $f_{\mathrm{score}}(C_i)$ with ABTC; (5) Verification vs persistent storage.

\textbf{ABTC (core contribution).} For each model--domain pair $(m,d)$, trust $\tau_m^{(d)} \sim \mathrm{Beta}(\alpha_m^{(d)}, \beta_m^{(d)})$; posterior mean $\hat{\tau}_m^{(d)} = \alpha/(\alpha+\beta)$. After each round, winner $m_w$ gets success, others failure, with decay $\lambda = 0.98$: $\alpha_m^{(d)} \leftarrow \lambda \alpha_m^{(d)} + \mathbb{1}[m=m_w]$, $\beta_m^{(d)} \leftarrow \lambda \beta_m^{(d)} + \mathbb{1}[m \neq m_w]$. Candidate score: $s_i = w_q s_i^{\mathrm{quality}} + w_a s_i^{\mathrm{agree}} + w_t \hat{\tau}_i$. If confidence $< \theta_{\min}$, synthesize from top-3.

\textbf{Model selection fitness:}
$\mathrm{fitness}(m,t) = 0.4\,\mathrm{CapMatch} + 0.4\,\mathrm{HistAcc} + 0.2(1-\hat{c}_m)$; SelectDiverseTop enforces provider diversity ($k_{\mathrm{models}}=3$, beam $k=3$).

% ----- Evaluation -----
\section*{Evaluation}
\textbf{Benchmark:} 500 queries, 5 domains (100 each): analytical reasoning, code generation, multi-step problem solving, knowledge retrieval, creative synthesis. 2--5 steps per query; context retrieval and cross-model validation.

\textbf{Quality:} GPT-4 as automated evaluator (LLM-as-judge); 50 responses human-annotated; Cohen's $\kappa = 0.74$, Pearson $r = 0.82$ vs human ($p < 0.001$).

\textbf{Baselines:} Sys-1 GAIOL; Sys-2 Direct API (single GPT-4); Sys-3 LangChain (ReAct + FAISS); Sys-4 OpenRouter (default routing); Sys-5 Multi-Wrapper (GPT-4 + Gemini Pro, confidence-based choice).

\textbf{Hardware:} Azure Standard\_D8s\_v3 (8 vCPU, 32\,GB RAM); 10 runs per config; means with 95\% CI.

% ----- Results -----
\section*{Results}
\begin{center}
\small
\resizebox{\columnwidth}{!}{%
\begin{tabular}{lccccc}
\toprule
Metric & S-1 & S-2 & S-3 & S-4 & S-5 \\
\midrule
Quality & $\mathbf{0.83}$ & 0.67 & 0.72 & 0.67 & 0.62 \\
Success (\%) & $\mathbf{95.2}$ & 92.0 & 94.5 & 93.8 & 90.5 \\
Latency (ms) & 450 & 400 & 800 & 420 & 600 \\
Overhead (ms) & 5 & 0 & 50 & 2 & 10 \\
\bottomrule
\end{tabular}%
}
\captionof{table}{Quality and performance (excerpt)}
\end{center}

GAIOL adds only 5\,ms orchestration overhead; per-query cost \$0.003 (comparable to single-model). ABTC ablation:

\begin{center}
\small
\resizebox{\columnwidth}{!}{%
\begin{tabular}{lccc}
\toprule
Domain & Equal & Tuned & ABTC \\
\midrule
Analytical & 0.79 & 0.81 & $\mathbf{0.86}$ \\
Code Gen. & 0.76 & 0.80 & $\mathbf{0.85}$ \\
Multi-step & 0.77 & 0.79 & $\mathbf{0.84}$ \\
Knowledge & 0.80 & 0.82 & $\mathbf{0.85}$ \\
Creative & 0.72 & 0.74 & $\mathbf{0.79}$ \\
\midrule
Overall & 0.77 & 0.79 & $\mathbf{0.83}$ \\
\bottomrule
\end{tabular}%
}
\captionof{table}{ABTC ablation ($p<0.01$ vs static)}
\end{center}

Largest gains: code generation (+9\,pp), creative synthesis (+7\,pp). Trust posteriors diverge by domain (e.g. GPT-4 $\hat{\tau}\approx 0.82$ analytical, $0.61$ creative; Gemini inverse), confirming domain-specific learning.

% ----- Conclusion -----
\section*{Conclusion}
GAIOL combines beam-search decomposition, strategy-based routing, and ABTC consensus in a single pipeline. ABTC replaces static weights with online-learned per-model, per-domain trust. On the 500-query benchmark, GAIOL achieves 0.83 quality, 95.2\% success, 5\,ms overhead; ABTC significantly outperforms uniform and hand-tuned consensus in all domains. Future work: federation and governance under multi-tenant/multi-cloud evaluation; standardized semantic protocols; adversarial robustness; GAIOL as backend for AutoGen/MetaGPT.

% ----- References -----
\section*{References}
\small
[1] P. Lewis et al., ``RAG for Knowledge-Intensive NLP,'' NeurIPS 2020.\\
[2] J. Wei et al., ``Chain-of-Thought Prompting,'' NeurIPS 2022.\\
[3] S. Yao et al., ``ReAct,'' ICLR 2023.\\
[4] Q. Wu et al., ``AutoGen,'' ICLR 2024.\\
[5] O. Chapelle \& L. Li, ``Thompson Sampling,'' NeurIPS 2011.\\
[6] H. Chase, ``LangChain,'' GitHub 2022.\\
[7] P. Liang et al., ``Holistic evaluation of language models,'' TMLR 2023.\\
[8] OpenAI, ``GPT-4 Technical Report,'' 2023.

\end{multicols}

\vspace{0.3cm}
\begin{center}
{\small Full manuscript and benchmark: see submitted paper to \textit{Engineering Applications of Artificial Intelligence}.}
\end{center}

\end{document}
